\scanpage{7}
% archive foliation: 6 (upper right and lower left)
en fait, \deleted{et} décrivant $\mathcal{M}(S)$.
Pratiquement, on n'a guère qu'à considérer dans \deleted{\unsure{une}} chaque question une ss-catégorie de type fini de $\mathcal{M}(S)$, soit $\mathcal C$. Elle est alors décrite par un \deleted{\unsure{proalgébrique}} groupe algébrique \inserted{réductif} $G$ sur $\mathbb{Q}$.
Question intéressante : comment sont reliés ces groupes pour des $\unsure{\xi|\zeta}$ différents (ils se déduisent tous les uns des autres par torsion\dots).
Déjà intéressant si $S$ est le spectre d'une ext.\ finie $k$ de $\mathbb{Q}$\dots!

Si $S$ est de car.\ $p>0$, on ne pourra en général trouver de foncteur fibre sur $\mathbb{Q}$ (contre-exemple de Serre, lié aux courbes elliptiques de Hasse nul).
Mais on aura un foncteur fibre sur $\mathbb{Q}_\ell$ pour $\forall$ pt géom de car.\ $\ne\ell$, d'où un foncteur fibre sur une extension finie \inserted{$K/$} de $\mathbb{Q}$ (si on se restreint à une $\mathcal C\subset \mathcal{M}(S)$, stable par \deleted{$\unsure{\otimes,\check{\phantom M}}$} $\otimes,\vee$, et à \unsure{engendrement} fini\dots).
Donc on trouve un groupe algébrique réductif $G_K$ sur $K$, qui permet de représenter la catégorie des $K'$-motifs pour $\forall$ corps de \unsure{coeff} $K'$ qui est une extension de $K$.

\margin{Structure $\varepsilon,\unsure{D}$, \unsure{graduations} $\unsure{\gamma_1},\unsure{\gamma_2}$
\inserted{$\unsure{\exp(D)=\lambda}$}
\deleted{\unsure{prolongement de} \illegible{2} $\unsure{C}$ \illegible{1}}
$\unsure{T}^{\unsure{*}}_{\mathrm{Hdg}}$.}
9) \uline{Relations avec le $\pi_1$, \unsure{profini}} :
\inserted{(A) $\unsure{G/G^0}\unsure{-|\longrightarrow}\unsure{\pi^1|\pi_1}$.}
Soit $F$ un foncteur fibre à coefficients dans $K$ sur $\unsure{\mathcal{M}(C)|\mathcal{M}(S)}$,
$\unsure{\xi|\zeta}$ un pt géométrique. On dit que $F$ et $\unsure{\xi|\zeta}$ sont associés si on a un \unsure{isom}
